\section{放射線の種類と材料への作用}
\label{theory:basicchem:radiation}
\(\alpha\)線はHe原子核、\(\beta^-\)線は電子、\(\beta^+\)線は陽電子、\(\gamma\)線は光子である。通常\(\alpha\)は強い電離作用と短い飛程を持つ。光子遮蔽の減衰は\(I=I_0e^{-\mu x}\)という単純モデルで表せ、\(\mu\)は材料と光子エネルギーに依存する\autocite{finish56_radiation}。同じ厚さと同じ面密度は別の比較である。照射は結晶欠陥や電子状態を変えて力学・電気特性を変え得る。

\section{構造と用途を対応させる化学物質}
\label{theory:basicchem:substances}
同素体は元素が同じでも結合・構造が違う。黒鉛は層状で軟らかく電気を伝え、ダイヤモンドの硬さ・通常の絶縁性と異なる。Haber--Boschはアンモニア合成、Solvayは炭酸ナトリウムの代表的製法である。
TiO2は光照射による表面反応を利用した抗菌・防汚等の光触媒材料として用いられる\autocite{finish56_titania}。性能は結晶相・複合化・光の波長等に依存する。鉛蓄電池の放電時は負極Pb、正極PbO2で、PbOという別酸化物と混同しない\autocite{openstax_batteries}。

\section{結合極性と分子極性を分ける}
\label{theory:basicchem:polarity}
極性結合を持つ分子でも、各結合の双極子が対称に配置されれば分子の双極子は0となる\autocite{finish56_polarity}。直線CO2、正四面体CH4・CCl4、平面三角形BF3は通常の形で相殺する。エーテルのO付近は折れ曲がり、分子の双極子が残る。結合を作る元素だけでなく、孤立電子対を含む分子形を確かめる。

\section{コロイドの粒子・界面・状態}
\label{theory:basicchem:colloids}
コロイドは分子溶液より大きい粒子等が分散する系で、粒子の界面と電荷が安定性を左右する。疎水コロイドに電解質を加えて反発を弱め、集合・沈殿させるのが凝析。小分子・イオンを膜で通して粒子と分けるのが透析である\autocite{openstax_colloids}。
光散乱で光路が見えるTyndall現象と、電場中で荷電粒子が移動する電気泳動は別の作用である。流動する分散系のゾルと、網目等ができ流動しにくいゲルを区別する。全コロイドが同じ量の電解質で凝析するわけではなく、表面・溶媒・イオン種等が関わる。
